\chapter{Studium przypadku I - Analiza poinfekcyjna i inżynieria detekcji (Próbka \#1)}\label{ch:studium_przypadku_1}

\section{Pozyskanie próbki, analiza wstępna i przebieg detonacji}\label{sec:cs1_triage}

\subsection{Metadane próbki i wstępna weryfikacja statyczna}\label{subsec:cs1_statyczna}
% TODO: Tabela z metadanymi Próbki #1 (nazwa rodziny, SHA-256 z makrem \hash{}, rozmiar, architektura PE, sekcje, entropia, podejrzane importy API oraz ciągi znaków z FLOSS).

\subsection{Kontrolowana egzekucja i akwizycja materiału dowodowego}\label{subsec:cs1_egzekucja}
% TODO: Opis uruchomienia Próbki #1 w środowisku Windows Victim, przebieg obserwacji oraz tabela zebranych plików dowodowych (.vmem, .evtx, .pcap) wraz z sumami SHA-256.

\section{Analiza śledcza pamięci operacyjnej RAM}\label{sec:cs1_analiza_ram}

\subsection{Identyfikacja procesów i anomalii w drzewie wykonania}\label{subsec:cs1_procesy_ram}
% TODO: Wyniki wtyczek Volatility 3 (windows.pslist, windows.pstree, windows.psscan, windows.cmdline) oraz zrzuty/listingi podejrzanych procesów.

\subsection{Detekcja wstrzykiwania kodu i analiza obszarów VAD}\label{subsec:cs1_iniekcja_vad}
% TODO: Identyfikacja obszarów PAGE_EXECUTE_READWRITE za pomocą windows.malfind / windows.vadinfo oraz analiza w MemProcFS.

\subsection{Artefakty sieciowe w pamięci i ekstrakcja zrzuconego ładunku (Payload)}\label{subsec:cs1_ekstrakcja_payload}
% TODO: Analiza połączeń w RAM (windows.netscan / netstat), ekstrakcja niezaszyfrowanego ładunku PE z pamięci procesu i jego analiza statyczna.

\section{Analiza telemetrii systemowej i ruchu sieciowego}\label{sec:cs1_analiza_logow}

\subsection{Automatyczny triage dzienników zdarzeń (Hayabusa i Chainsaw)}\label{subsec:cs1_triage_evtx}
% TODO: Wyniki szybkiego skanowania logów .evtx narzędziami Hayabusa i Chainsaw, zestawienie wykrytych alertów według krytyczności.

\subsection{Rekonstrukcja aktywności procesowej, plikowej i rejestrowej (Sysmon)}\label{subsec:cs1_sysmon_analiza}
% TODO: Szczegółowa analiza zdarzeń Sysmon (Event ID 1, 7, 8, 10, 11, 12/13) - parametry CLI, iniekcja, persystencja w rejestrze lub harmonogramie zadań.

\subsection{Analiza przechwyconej komunikacji C2 (INetSim / PCAP)}\label{subsec:cs1_analiza_sieci}
% TODO: Analiza zapytań DNS (Sysmon Event ID 22 + PCAP) oraz przechwyconych sesji HTTP/HTTPS/TCP na bramie REMnux.

\section{Korelacja artefaktów i rekonstrukcja łańcucha ataku}\label{sec:cs1_korelacja}
% TODO: Chronologiczna tabela korelacji (Timestamp UTC, źródło: RAM / Sysmon / PCAP, PID/ProcessGuid, opis zdarzenia) oraz tabela mapowania TTP na matrycę MITRE ATT&CK.

\section{Inżynieria detekcji i walidacja reguł}\label{sec:cs1_inzynieria_detekcji}

\subsection{Projekt i implementacja reguły YARA}\label{subsec:cs1_regula_yara}
% TODO: Listing autorskiej reguły YARA (style=code) dla wyekstrahowanego ładunku z pamięci RAM wraz z uzasadnieniem doboru ciągów i warunków.

\subsection{Projekt i implementacja reguł Sigma}\label{subsec:cs1_regula_sigma}
% TODO: Listing autorskiej reguły Sigma (style=code) dla zachowań zarejestrowanych w dziennikach Sysmon wraz z opisem logiki detekcji.

\subsection{Weryfikacja skuteczności i testy fałszywych alarmów (FP/TP)}\label{subsec:cs1_walidacja}
% TODO: Wyniki walidacji reguł YARA i Sigma na zainfekowanym materiale dowodowym (TP) oraz na czystej migawce referencyjnej Clean_State (FP = 0).

\section{Zestawienie wskaźników kompromitacji (IoC) oraz aktywności (IoA)}\label{sec:cs1_ioc_ioa}
% TODO: Zbiorcza tabela wskaźników IoC (hashe SHA-256, ścieżki plików, klucze rejestru, mutexy, domeny/IP C2) oraz wskaźników behawioralnych IoA dla Próbki #1.

